(1) 用快速度$\eta$表示四速度，$U^\mu=(c\cosh\eta,c\sinh\eta)$；恒定固有加速度要求$d\eta/d\tau=a/c$，初值$\eta(0)=0$。积分$dt/d\tau=\cosh(a\tau/c)$、$dx/d\tau=c\sinh(a\tau/c)$，得到$t(\tau)=(c/a)\sinh(a\tau/c)$、$x(\tau)=(c^2/a)[\cosh(a\tau/c)-1]$，实验室速度$v=c\tanh(a\tau/c)$。
(2) 从上述坐标消去$\tau$，得$(x+c^2/a)^2-c^2t^2=(c^2/a)^2$，为右支双曲线。四加速度$dU^\mu/d\tau=(a\sinh\eta,a\cosh\eta)$，闵氏范数$A^\mu A_\mu=-a^2$（取度规$+,-$），因此固有加速度大小始终为$a$，即使实验室坐标加速度$\ddot x=a/\gamma^3$随速度增大而趋零。
(3) 若实验室原点在$t_e\ge0$发出沿$+x$光束，其轨迹为$x_{\gamma}=c(t-t_e)$。火箭轨迹满足$ct-x=(c^2/a)(1-e^{-a\tau/c})$，右边随$\tau$从0单调上升但上界为$c^2/a$。两世界线相交要求$ct_e<c^2/a$，即$t_e<c/a$；在$t_e=c/a$时只在无穷远极限相遇，之后发射的光永远追不上火箭。此为加速观察者的事件视界效应。
(4) 同向追赶火箭的实验室光子四波矢$(\omega/c,\omega/c)$，火箭测得频率$\omega'=\gamma\omega(1-v/c)=\omega(\cosh\eta-\sinh\eta)=\omega e^{-\eta}$。所以$\omega'(\tau)=\omega\exp[-a\tau/c]$。越接近视界的追赶光子在火箭上越红移；这里假定实验室光源频率$\omega$不变且忽略引力曲率。
